El presente capítulo establece las reglas generales, políticas técnicas y procedimientos básicos aplicables a todas las actividades de mantenimiento, inspección, reparación y control técnico de los Sistemas de Aeronaves No Tripuladas (UAS) de SKYFILMS, con el fin de garantizar que dichas labores se ejecuten conforme a los estándares definidos por el fabricante y a los Reglamentos Aeronáuticos de Colombia (RAC 100).\\

Estas reglas son de cumplimiento obligatorio para todo el personal técnico y administrativo involucrado en los procesos de mantenimiento, asegurando la preservación de la aeronavegabilidad, la seguridad operacional y la trazabilidad de la información técnica.

\section{Principios Generales}

\begin{enumerate}
    \item Todas las actividades de mantenimiento deberán realizarse conforme a lo establecido en este Manual de Control de Mantenimiento (MCM) y en los manuales técnicos y boletines del fabricante de cada UAS.

    \item Ningún mantenimiento podrá realizarse fuera de las instalaciones aprobadas o por personal no autorizado por SKYFILMS.

    \item Todo trabajo técnico deberá documentarse en los formatos oficiales aprobados, garantizando la trazabilidad completa de cada intervención.

    \item Se prohíbe expresamente el uso de herramientas, equipos o repuestos no aprobados por el fabricante o sin registro de control.

    \item El mantenimiento debe realizarse bajo condiciones de orden, limpieza y seguridad, evitando contaminación de componentes, sobrecargas eléctricas o exposición a humedad.

    \item Los reportes de discrepancias detectadas durante la operación deberán comunicarse inmediatamente al área técnica antes de la próxima operación del UAS.
\end{enumerate}

\section{Tipos de Mantenimiento Aplicables}
SKYFILMS reconoce los siguientes tipos de mantenimiento aplicables a su flota de UAS:

\begin{itemize}
    \item \textbf{Mantenimiento menor:} Comprende las tareas rutinarias, preventivas y de baja complejidad ejecutadas para conservar el equipo en óptimas condiciones operativas. Incluye actividades como inspecciones visuales calendarizadas, limpieza superficial, ajustes mecánicos básicos y la sustitución de componentes de desgaste normal o de fácil acceso (por ejemplo, reemplazo de hélices o protectores). Estas tareas no requieren un desensamblaje profundo ni pruebas de verificación funcional complejas para autorizar el retorno al servicio.

    \item \textbf{Mantenimiento mayor:} Involucra intervenciones de alta complejidad técnica, estructural o electrónica que requieren un nivel superior de especialización. Este renglón incluye la reparación o reemplazo de componentes críticos (tales como motores, controladores de vuelo, ESCs o estructuras del brazo/fuselaje), calibraciones profundas de sistemas de navegación, y actualizaciones de firmware que alteren los parámetros de vuelo. Este tipo de mantenimiento siempre requiere una prueba de verificación funcional posterior para autorizar la liberación de la aeronave.

    \item \textbf{Mantenimiento correctivo:} Comprende las acciones no programadas y necesarias para corregir fallas, discrepancias técnicas o averías detectadas durante inspecciones previas o durante la operación de vuelo. Se activa tras la notificación de una anomalía o evento inusual (ej. aterrizajes fuertes, impactos, exposición a humedad o mensajes de error en telemetría) y su objetivo exclusivo es restituir la aeronave a su condición de aeronavegabilidad certificada y segura.

\end{itemize}

En ningún caso el personal de SKYFILMS ejecutará tareas que excedan las competencias o procedimientos autorizados por el fabricante o por este manual.

\section{Procedimientos Generales de Ejecución}

\begin{itemize}
    \item Antes de iniciar cualquier tarea de mantenimiento, el técnico responsable deberá verificar la aeronavegabilidad del equipo, la vigencia de los registros y la disponibilidad de las herramientas necesarias.

    \item Toda intervención debe realizarse con el sistema de energía desconectado y con las baterías retiradas, salvo en los casos en que el procedimiento requiera pruebas eléctricas bajo control supervisado.

    \item El área de trabajo debe mantenerse ordenada y libre de elementos ajenos al mantenimiento.

    \item Los componentes sustituidos deben registrarse en el Formato de Reemplazo de Componentes y etiquetarse para su disposición o análisis posterior.

    \item Las tareas finalizadas deben ser verificadas y firmadas y revisadas por el Técnico UAS antes de liberar la aeronave.
\end{itemize}

\section{Modificaciones y Alteraciones}
\begin{itemize}
    \item Toda modificación o alteración de la configuración o de las características técnicas que afecte las capacidades o las características de vuelo de una aeronave debe estar aprobada expresamente por su fabricante (RAC 100.415).
    \item En el Mavic 4 Pro solo se instalan repuestos y accesorios originales o aprobados por DJI. En el Pavo20 Pro II, los cambios de motores, hélices, batería, receptor, GPS o iluminación solo se hacen con componentes compatibles con el fabricante (BetaFPV) y se registran en el Libro de Mantenimiento, con una prueba funcional antes de volar.
    \item Un cambio que altere el peso máximo, el sistema de enlace C2 o los sistemas de fábrica se informa al Jefe de Pilotos y al Gerente SMS, y se actualizan este Manual, el Manual de Operaciones y las OpSpecs antes de operar.
    \item Los sistemas de fábrica (luces, sensores y demás partes originales) deben estar funcionales; una aeronave con un sistema de fábrica inoperativo no se libera para vuelo.
\end{itemize}

\section{Control de Documentación Técnica}

\begin{itemize}
    \item Todos los procedimientos de mantenimiento deben estar respaldados por documentos técnicos válidos, actualizados y disponibles en la biblioteca técnica de SKYFILMS.

    \item El Técnico de UAS es responsable de mantener actualizada la documentación del fabricante y de notificar cualquier cambio que pueda afectar los procedimientos.

    \item Las enmiendas o modificaciones en manuales técnicos sólo serán válidas una vez aprobadas por el Técnico de UAS y registradas en el Control de Documentos Técnicos.
\end{itemize}

\section{Registro de Mantenimiento}

\begin{itemize}
    \item Cada tarea ejecutada deberá ser registrada en el \textbf{Libro de Mantenimiento UAS}, indicando fecha, descripción del trabajo, técnico responsable, horas de operación y firma de aprobación.

    \item Los registros deberán conservarse de manera ordenada y accesible, por un período mínimo de cinco (5) años, igual que los registros del Manual de Operaciones, o el que exija la Aerocivil.

    \item El uso de registros electrónicos está permitido, siempre que garantice la integridad, autenticidad y disponibilidad de la información.
\end{itemize}

\section{Seguridad y Condiciones de Trabajo}
Durante la ejecución del mantenimiento, el personal técnico deberá cumplir las normas internas de seguridad y salud en el trabajo (SST) y las siguientes disposiciones específicas:

\begin{itemize}
    \item Uso obligatorio de elementos de protección personal (EPP) según el tipo de tarea.

    \item Prohibición de fumar, consumir alimentos o bebidas en el área técnica.

    \item Eliminación de residuos y componentes defectuosos conforme a la normativa ambiental vigente.

    \item Inmediata interrupción del mantenimiento ante cualquier condición insegura detectada.
\end{itemize}

\section{Control de Liberación de Aeronaves}

\begin{itemize}
    \item Ningún UAS podrá ser liberado para vuelo sin la autorización del Tecnico UAS, quien verificará que se haya completado toda la documentación técnica y que la aeronave cumpla con los requisitos de aeronavegabilidad.

    \item La firma de liberación constituye la garantía de que el UAS se encuentra en condiciones técnicas seguras y conforme con los procedimientos establecidos.

    \item En caso de detectarse discrepancias, la aeronave permanecerá inoperativa hasta la corrección total de la anomalía.
\end{itemize}

\section{Responsabilidades}

\begin{itemize}
    \item El \textbf{Técnico UAS} es responsable de la aplicación y cumplimiento de las reglas contenidas en este capítulo, así como de la ejecución correcta de los procedimientos y el diligenciamiento de registros.

    \item El \textbf{Jefe de Pilotos} deberá supervisar los procedimientos establecidos y reportar cualquier desviación o no conformidad.
\end{itemize}

\section{Disposiciones Finales}
Cualquier desviación o incumplimiento de las reglas establecidas en este capítulo será considerada una no \textbf{conformidad técnica} y deberá ser reportada al Gerente General y al Gerente de Seguridad Operacional, aplicando los procedimientos de reporte y acción correctiva establecidos por el Sistema de Gestión de Seguridad Operacional (SMS) de SKYFILMS.